% ================================================================
% ==== FILE: content/cycles/cycles_k1.tex
% ================================================================

\subsubsection{Zyklen auf \texorpdfstring{$K_1$}{K_1}}
\label{sec:cycles-k1}

Ebene$K_1$ist das erste echte Kontinuum in der Ontologie von
Fortsetzung.
Es wird durch eine einzige durchgehende Achse definiert
\[
  A_1(x)=x, \qquad x\in I=(a,b),
\]
und ein nicht leerer Zustandsraum
\[
  \Omega(K_1)
  = C^0\!\bigl(T, H^1(I,V)\bigr)
    \cap C^1\!\bigl(T, L^2(I,V)\bigr),
\]
wie in der formalen Konstruktion von festgelegt$K_1$.

Time $\tau(K_1)$ is first definable here, enabling the introduction of
Zyklen.
daher$K_1$ist das minimale Niveau, auf dem Zyklen existieren.

\subsubsection{Definition von Zyklen auf \texorpdfstring{$K_1$}{K_1}}

Let $s(t)\in\Omega(K_1)$ denote the state of the continuum.  
Ein Zyklus ist eine Abbildung
\[
  C:[0,\tau]\to\Omega(K_1), \qquad
  C(0)=C(\tau),
\]
vom Evolutionsoperator generiert
\[
  \frac{dA_1}{dt} = F_1(A_1,P_1,J_1,\Theta_1),
\]
subject to the admissibility conditions of $\Omega(K_1)$.

Weil$K_1$hat nur eine einzige Achse und darüber hinaus keine internen Schwellenwerte
$\Theta_1$, the dynamics is one-dimensional and cannot create
nichttriviale topologische oder strukturelle Zykluskomplexe.
Alle Zyklen laufen$K_1$sind daher in der Thermodynamik „trivial“ und
Struktureller Sinn.

\subsubsection{Eigenschaften von Zyklen auf \texorpdfstring{$K_1$}{K_1}}

Zyklen auf dieser Ebene weisen mehrere charakteristische Merkmale auf:

\begin{itemize}
  \item \textbf{Topologie:}
Der Zyklus lebt von einer 1D-Mannigfaltigkeit; Keine Schleifenverschachtelung oder Schleife
Komplexe existieren.

  \item \textbf{Entartung:}
Der Zyklus ist typischerweise eine minimale Schwingung oder ein geschlossener Zyklus
1-Parameter-Trajektorie des Feldes$A_1(x,t)$.

  \item \textbf{Keine kritischen Schwellenwerte:}
        The threshold vector $\Theta_1$ contains only the basic
Zulässigkeitsgrenze, die sich aus der Konstruktion von ergibt$K_1$.
Es sind keine Phasenübergänge damit verbunden$C$.

  \item \textbf{Keine strukturelle Spannung:}
Die strukturelle Spannung
$T_1(t)$ist definiert, kann aber keine komplexen Bifurkationen erzeugen oder
Instabilitäten, weil$K_1$lässt keine unvereinbaren Unterschiede zu.

  \item \textbf{Grenzverhalten:}
        The boundary $\partial\Omega(K_1)$ is trivial; cycles cannot
sinnvoll annähern oder überschreiten.

  \item \textbf{Energiekonsistenz:}
Die Aktion funktional
        \[
          S[A_1] = \int_0^\tau E(A_1,P_1,J_1)\, dt
        \]
bestimmt die dynamische Durchführbarkeit eines Zyklus, tut dies jedoch nicht
Erstellen Sie Einschränkungen höherer Ordnung.
\end{itemize}

\subsubsection{Metriken für \texorpdfstring{$K_1$}{K_1} Zyklen}

Aus dem in Run~9 entwickelten universellen Zyklusformalismus folgen Zyklen$K_1$
haben:

\[
  L(C) = \int_0^\tau \| \dot{A}_1(t) \|\, dt,
\]
\[
  C_{\mathrm{eff}} =
    \frac{1}{L(C)} \int_0^\tau J_1(t)\cdot dA_1(t),
\]
und eine triviale Stabilität
\[
  S(C) = \min_{t\in[0,\tau]} d(\Omega(K_1),\partial\Omega(K_1)).
\]

Da die Grenze trivial und unerreichbar ist,$S(C)$ist immer
maximal für zulässige Zustände:
\[
  S(C) = S_{\max}.
\]

\subsubsection{Rolle von \texorpdfstring{$K_1$}{K_1} Zyklen in der Hierarchie}

Die Existenz von Zyklen auf$K_1$ist für die höheren Ebenen unerlässlich:

\begin{itemize}
  \item Sie legen die erste zeitliche Struktur ungleich Null fest.
  \item Sie dienen als „Samen“-Zyklen, aus denen hervorgeht$K_2$
erbt Periodizität und dynamische Wiederkehr.
  \item Sie bilden die Grundlage für die nicht triviale Struktur
        cycles (e.g.\ percolation cycles, metabolic cycles, spike
Zyklen) entstehen bei$K_2$–$K_5$.
\end{itemize}

Also zwar Zyklen weiter$K_1$sind dynamisch einfach, das sind sie
strukturell unabdingbar für die Kontinuität der Hierarchie.

\subsubsection{Zusammenfassung}

$K_1$unterstützt die einfachsten möglichen Kreisläufe – geschlossene Trajektorien von a
1D-Kontinuum ohne Kritikalität, ohne Verzweigungen und ohne interne
Struktur.
Sie markieren die Geburt der Zeit und die dynamische Wiederkehr, die alles ermöglicht
Zyklen höherer Ordnung in der Ontologie von Continua.
